\section{Implementation}
\label{sec:impl}

The methods of Sections~\ref{sec:wf-topdown}, \ref{sec:dir-opt} and~\ref{sec:variants}, and their proofs, are stated in the model of Section~\ref{sec:bg-model}, which speaks of loads, stores, CASes and fences of words rather than of any programming language. An implementation keeps the guarantees of the proofs if it realizes these operations as the model describes them: each atomic, in the order that the model keeps, and each in a bounded number of steps of its thread. How to do so depends on the language, the compiler and the processor. This section describes our implementations of the three methods, which are written in C++ and built with GCC for x86-64 (Section~\ref{sec:exp-setup}).

In C++, two accesses to the same location by different threads, at least one of them a write, form a data race unless both are atomic, and a program with a data race has undefined behavior~\cite{boehm_cpp}. A compiler may therefore assume that no other thread changes a location while a thread accesses it with ordinary loads and stores, and transform these accesses in ways that are harmless for one thread but not for several: it may split an access into smaller ones, or load a location again instead of keeping a value that it loaded, so that two uses of one read in the source see different values~\cite{boehm_benign}. Our methods share locations in exactly this way: helpers read buckets, progress fields and flags while their owners write them, and a thread that has fallen behind writes distances, codes and stamps that other threads are reading.

Our implementations therefore make every access to a location that another thread may access at the same time, with one of the two a write, an atomic one. Most are \emph{relaxed} loads and stores, which are atomic but order nothing: those of distances, codes and flags $\mathit{vis}$; of the sizes, arrays and positions of buckets; of the progress fields $\mathit{last}$, $\mathit{done}$, $\mathit{codesLast}$ and $\mathit{codesDone}$ and the edge counts of entries; of stamps; and, in the variant, of the tags of buckets. We make them with \texttt{std::atomic\_ref}, part of C++ since C++20, which applies atomic operations to an ordinary object. The arrays thus stay ordinary arrays, and the code that sets them up and resets them between searches, when no other thread accesses them, stays ordinary code. A relaxed load or store is indivisible, and the compiler neither loads the location again in place of a value already loaded, nor stores to it where the source does not.

Ordinary accesses remain for memory that only one thread uses during a search, such as the capacity of a bucket; for memory that no thread writes during a search, such as the graph; and for memory that a thread writes before publishing it with a CAS, and that other threads read only after reading what the CAS wrote: a node, which its thread sets up before linking it (Lemma~\ref{lem:list}), and a proposal, which its thread writes before claiming the level (Lemma~\ref{lem:plans}). The fields $\mathit{next}$ and $\mathit{planner}$, and the word of a node in the variant, are \texttt{std::atomic} objects, whose loads and CASes are sequentially consistent. In C++, a load that reads the value written by such a CAS orders the accesses that came before the CAS before those that come after the load, so these ordinary accesses are not data races either.

In the variant, a thread that reads a reused bucket can read positions of its array that no thread has written in the search (Section~\ref{sec:mem-reuse}). The model lets it read them, since every word holds some value, but C++ gives memory that a program has never written an indeterminate value, which a program may not read. The arenas of the variant therefore take their memory from the system zeroed, with \texttt{calloc}, so that every position holds 0 or a value written in an earlier search, which \textsc{Valid} discards unless it is a vertex of the reader's level.

The distances are final once the first call of a search returns, but calls that are still running may write them again, with the values already there (Section~\ref{sec:bg-method}). A consumer that reads them before every call has returned must therefore read them with atomic loads too, or first copy them with atomic loads into memory of its own; it finds the final distances either way.

Relaxed accesses can be reordered: the compiler may move them past each other, and a processor whose memory model is weaker than TSO may make them take effect out of order. On x86-64, GCC compiles every relaxed load and store to a single \texttt{mov} instruction, which TSO keeps in program order with the other accesses of its thread, except a store and a later load, so only the compiler can break the orders that the model needs. We keep that order with \emph{compiler barriers}, empty \texttt{asm} statements that clobber memory, which emit no instruction, and across which GCC moves no memory access. We put one wherever the algorithms require one access of a thread to come before another: between every write that lets other threads skip work and the work it covers (Section~\ref{sec:td-order}); between appending a vertex, writing its code, and writing its distance; in \textsc{Append}, between writing a vertex and the new size, and between copying an array and publishing it (Algorithm~\ref{alg:bucket}); between reading the size of a bucket and its array; and, in the variant, between resetting the progress of a thread, emptying its bucket, and tagging it, and between reading the tag of a bucket and the size or the progress that the tag guards (Algorithms~\ref{alg:mem} to~\ref{alg:mem-codes}). With these barriers, every thread makes its loads and stores in the order that the algorithms list them, and TSO keeps that order as Section~\ref{sec:bg-model} describes.

This way of keeping the order relies on GCC and on TSO. Portable C++ would put the order into the accesses themselves, with a release store for every store that must become visible after the earlier loads and stores of its thread, and an acquire load for every load that must take effect before the later loads and stores of its thread~\cite{boehm_cpp}. On x86-64, GCC compiles release stores and acquire loads to the same \texttt{mov} instructions as relaxed ones, so on our machine the two forms make the same instruction for every access. Processors with weaker memory models, such as ARM and POWER, need instructions of their own for release and acquire, and there the barriers alone would not keep the order.

Every CAS of the methods is a \texttt{compare\_exchange\_strong} with the default, sequentially consistent order, and the fence of Algorithm~\ref{alg:wf-td} (line~\ref{ln:td-fence}) is an \texttt{atomic\_thread\_fence} with the same order. GCC compiles the CAS to a \texttt{lock cmpxchg} instruction, and the fence to a \texttt{lock or} of zero into the top of the stack, which changes no data. A locked instruction takes effect only once the earlier stores of its thread are visible~\cite{intelmanual}, and GCC moves no memory access across a sequentially consistent operation, so both behave as Section~\ref{sec:bg-model} assumes. The fence costs one instruction per level.

The proof of wait-freedom counts every load, store, CAS and fence as one step (Section~\ref{sec:bg-model}), so it carries over to an implementation only if each of them finishes in a bounded number of steps of its own thread. C++ implements an atomic operation that the processor cannot perform directly with a lock, which a delayed thread could hold; every atomic type that we use, of one, four or eight bytes, is always lock-free on x86-64 (\texttt{is\_always\_lock\_free}), and each of our atomic operations is a single instruction. A thread's arena hands out blocks by moving a pointer through chunks of at least 4~MiB, which it keeps from one search to the next, and asks the system allocator for a new chunk only when none of its chunks has room left. These calls, which can take a lock, are the only ones outside the assumption of Section~\ref{sec:bg-model}, and a thread makes them only in a search that needs more memory than its arena already holds.
